%% v3_dalmau.tex — environment specification
% !TEX program = pdflatex

\documentclass[11pt, a4paper]{article}
\usepackage{amsmath, amssymb}
\usepackage{booktabs}
\usepackage{geometry}
\usepackage{microtype}
\usepackage{parskip}
\usepackage[hidelinks]{hyperref}
\geometry{margin=2.5cm}

\title{\textbf{v3\_dalmau: Environment Specification}}
\author{}
\date{}

\begin{document}
\maketitle

% ─────────────────────────────────────────────────────────────────────────────
\section{Overview}
% ─────────────────────────────────────────────────────────────────────────────

The \texttt{v3\_dalmau} environment is a single-sector air traffic conflict
resolution environment in which a reinforcement learning agent issues discrete
heading instructions to maintain separation between en-route aircraft.
Each episode places between ten and fifteen Airbus A320 aircraft in a randomly
shaped convex sector at FL350.
Aircraft enter from the sector boundary, transit toward a destination beyond
the opposite boundary, and are replaced when they exit, keeping traffic density
approximately constant.
The agent observes the state of the focus aircraft and its four most urgent
neighbours and selects one of eight discrete actions per step.

The action space is structured around the ATC resolution pattern of turn,
hold, and return.
Six \textit{turn instructions} deflect the focus aircraft's commanded heading
by $\delta \in \{{\pm}30^\circ,\,{\pm}45^\circ,\,{\pm}60^\circ\}$ from the
destination bearing.
A \textit{hold} action maintains the current commanded heading without
additional cost, allowing the aircraft to execute a turn over multiple steps.
A \textit{fly-direct} action returns the commanded heading to the destination
bearing, eliminating accumulated drift once the conflict has resolved.

The reward signal is purely negative: a large penalty is applied when a
separation violation occurs, a graduated conflict penalty scales with the
imminence and predicted miss distance of approaching pairs, a drift penalty
accumulates while the aircraft is held off its planned route, and a workload
penalty is charged once per turn instruction.
The combination of drift and workload terms creates a sustained incentive to
resolve conflicts with as few instructions and as little deviation as possible.

The focus aircraft, the one to which the instruction is applied, is selected
each step by an urgency-based mechanism that prioritises aircraft in or near
separation violations and locks onto a single aircraft for the duration of a
resolution sequence.

% ─────────────────────────────────────────────────────────────────────────────
\section{Scenario}
\label{sec:scenario}
% ─────────────────────────────────────────────────────────────────────────────

Episodes are simulated using the BlueSky air traffic simulator at a physics
timestep of $\Delta t = 0.5$\,s.
At the start of each episode, $N \sim \mathcal{U}\{10,\ldots,15\}$ Airbus A320
aircraft are initialised at FL350 and cruise at a fixed speed of
$M = 0.78$ ($\approx 450$\,kt TAS) throughout.
The sector boundary is a random convex polygon with five to seven vertices,
scaled to an area of $N \times \rho$\,km$^2$, where
$\rho \sim \mathcal{U}(5\,000,\,15\,000)$\,km$^2$/aircraft.
Polygons with an isoperimetric ratio below 0.65 are rejected to avoid
pathologically elongated sectors.
Each aircraft is spawned at a distinct arc of the sector perimeter with
uniform jitter, assigned a heading toward a reference point on the opposite
arc, and given a destination two sector diameters beyond the boundary.
A minimum spawn separation of 15\,NM is enforced.
Aircraft that exit the sector are replaced within one step, keeping traffic
density approximately constant throughout each episode.
Episode length corresponds to approximately 60 minutes of simulated time,
derived from the sector diameter and nominal cruise speed.

Each RL step advances the simulation by $10 \times 0.5\,\text{s} = 5\,\text{s}$.
The selected action is applied at the start of the step and one instruction
is issued to the \textit{focus aircraft}, selected as described in
Section~\ref{sec:focus}.
Separation is checked at the end of the final sub-step.

% ─────────────────────────────────────────────────────────────────────────────
\section{Focus Selection}
\label{sec:focus}
% ─────────────────────────────────────────────────────────────────────────────

Because only one instruction is issued per step, a principled selection
mechanism is required to determine which aircraft receives it.
An urgency matrix $\mathbf{U} \in \mathbb{R}^{N \times N}$ is recomputed
after every step from current positions and velocities.
The entry $u_{ij}$ is defined as

\begin{equation}
u_{ij} =
\begin{cases}
1 + 9\!\left(1 - \dfrac{d_{ij}}{d_\text{sep}}\right)
  & d_{ij} < d_\text{sep} \\[8pt]
\dfrac{T_\text{warn} - t_\text{CPA}^{ij}}{T_\text{warn}}
  & d_\text{CPA}^{ij} < d_\text{sep},\quad
    0 \leq t_\text{CPA}^{ij} \leq T_\text{look} \\[6pt]
0 & \text{otherwise,}
\end{cases}
\label{eq:urgency}
\end{equation}

with $d_\text{sep} = 5$\,NM, $T_\text{warn} = 600$\,s, and
$T_\text{look} = 900$\,s.
The active-violation branch scales linearly from $u_{ij} = 1$ at the
separation boundary to $u_{ij} = 10$ at zero distance.
The predicted-conflict branch ramps linearly from 0 at $T_\text{warn}$ to 1
at $t_\text{CPA} = 0$; both branches are continuous at $d_{ij} = d_\text{sep}$.
The focus aircraft is selected as $f = \arg\max_i \max_{j \neq i}\, u_{ij}$.
When no conflicts are active, focus falls back to the aircraft with the
largest commanded-heading deviation from its destination, subject to a
hysteresis margin of $0.05$ to prevent excessive switching.

\paragraph{Commitment.}
Conflict resolution typically requires several consecutive instructions
to the same aircraft.
The focus is therefore locked on $f$ while $u_f > 0$ or fewer than
$n_\text{clear} = 5$ steps have elapsed since the urgency last cleared.
An emergency override releases the lock immediately when any other aircraft
reaches $u_\text{emerg} = 0.8$, corresponding to approximately 120\,s
to predicted conflict.

% ─────────────────────────────────────────────────────────────────────────────
\section{Action Space}
\label{sec:actions}
% ─────────────────────────────────────────────────────────────────────────────

The action space is discrete with eight elements, summarised in
Table~\ref{tab:actions}.
The eight actions fall into three groups.

\paragraph{Turn instructions (actions 0--2 and 4--6).}
Six actions set the commanded heading of the focus aircraft as a one-shot
offset from the current destination bearing:
\begin{equation}
    \psi_\text{cmd} \leftarrow \bigl(\psi_\text{dest}(f) + \delta\bigr)\bmod 360^\circ,
    \qquad
    \delta \in \{-60^\circ,\,-45^\circ,\,-30^\circ,\,+30^\circ,\,+45^\circ,\,+60^\circ\}.
\end{equation}
The offset is computed relative to the live direct-to-destination bearing at
the moment the action is applied, so the commanded deviation is always
expressed in terms of the current route.
A workload cost $c$ is charged once when the instruction is issued; larger
deviations carry higher cost.

\paragraph{Hold (action 3).}
Hold re-issues $\psi_\text{cmd}$ unchanged, allowing the aircraft to continue
executing a previously commanded turn without incurring an additional
instruction cost.
It is the only action with zero cost and serves as the natural
\textit{no-intervention} baseline: if hold is applied every step, all
aircraft follow their commanded headings without further modification.

\paragraph{Fly direct (action 7).}
Fly direct resets $\psi_\text{cmd} \leftarrow \psi_\text{dest}(f)$,
returning the aircraft to its planned route in a single step.
The cost scales with the commanded deviation being corrected, so that
fly direct is effectively free when the aircraft is already on track
and expensive when a large turn is being aborted prematurely.

\medskip
The three groups encode the intended ATC resolution pattern: issue one turn
instruction to create lateral separation, hold while the aircraft manoeuvres
and the conflict geometry resolves, then fly direct to eliminate accumulated
route deviation.

\begin{table}[h]
\centering
\caption{Action space and associated workload costs.}
\label{tab:actions}
\smallskip
\begin{tabular}{clc}
\toprule
\textbf{Index} & \textbf{Action} & \textbf{Cost $c$} \\
\midrule
0 & $\delta = -60^\circ$ & 1.5 \\
1 & $\delta = -45^\circ$ & 1.0 \\
2 & $\delta = -30^\circ$ & 0.5 \\
3 & Hold ($\psi_\text{cmd}$ unchanged) & 0.0 \\
4 & $\delta = +30^\circ$ & 0.5 \\
5 & $\delta = +45^\circ$ & 1.0 \\
6 & $\delta = +60^\circ$ & 1.5 \\
7 & Fly direct ($\psi_\text{cmd} \leftarrow \psi_\text{dest}$) & $|\Delta\psi_\text{cmd}|/60^\circ \in [0,1]$ \\
\bottomrule
\end{tabular}
\end{table}

% ─────────────────────────────────────────────────────────────────────────────
\section{Observation Space}
\label{sec:obs}
% ─────────────────────────────────────────────────────────────────────────────

The observation vector $\mathbf{o} \in \mathbb{R}^{23}$ is constructed
ego-centrically from the focus aircraft $f$.
It comprises three ownship features and four intruder slots of five features
each.
All positional and velocity quantities are expressed in $f$'s heading frame,
with the forward axis aligned with $\psi_\text{act}$ and the lateral axis
$90^\circ$ to the right.

\begin{table}[h]
\centering
\caption{Observation vector ($\mathbf{o} \in \mathbb{R}^{23}$).
$d_\text{warn} = T_\text{warn}\,v_\text{nom} = 75$\,NM.
CPA coordinates are normalised by $d_\text{sep}$, the correct scale for
conflict geometry.}
\label{tab:obs}
\smallskip
\begin{tabular}{clll}
\toprule
\textbf{Index} & \textbf{Symbol} & \textbf{Definition} & \textbf{Range} \\
\midrule
\multicolumn{4}{l}{\textit{Ownship state}} \\[2pt]
0 & $\sin\Delta\psi$ & $\sin(\psi_\text{dest} - \psi_\text{cmd})$ & $[-1,1]$ \\
1 & $\cos\Delta\psi$ & $\cos(\psi_\text{dest} - \psi_\text{cmd})$ & $[-1,1]$ \\
2 & $p_\text{turn}$  & $\operatorname{clip}\bigl((\psi_\text{cmd}-\psi_\text{act})/60^\circ,\,-1,1\bigr)$ & $[-1,1]$ \\[6pt]
\multicolumn{4}{l}{\textit{Per intruder $j$ (slots 1--4, offset $+k$ for slot $k$)}} \\[2pt]
$+0$ & $\tilde{x}_{fj}$            & lateral relative position $/ \, d_\text{warn}$            & $[-1,1]$ \\
$+1$ & $\tilde{y}_{fj}$            & forward relative position $/ \, d_\text{warn}$            & $[-1,1]$ \\
$+2$ & $\tilde{x}^\text{CPA}_{fj}$ & lateral CPA offset $/ \, d_\text{sep}$                    & $[-1,1]$ \\
$+3$ & $\tilde{y}^\text{CPA}_{fj}$ & forward CPA offset $/ \, d_\text{sep}$                    & $[-1,1]$ \\
$+4$ & $\hat{t}_\text{CPA}$        & $\operatorname{clip}(t_\text{CPA}/T_\text{warn},\,0,\,1)$  & $[0,1]$ \\
\bottomrule
\end{tabular}
\end{table}

The heading error is encoded as the pair $(\sin\Delta\psi,\,\cos\Delta\psi)$
rather than a scalar angle, which avoids the discontinuity at $\pm 180^\circ$.
The error is measured against the \textit{commanded} heading $\psi_\text{cmd}$
rather than the actual heading, consistent with the drift penalty in the
reward function.
Turn progress $p_\text{turn}$ captures the execution lag: it is zero when
the aircraft has completed its commanded turn and non-zero while the physical
heading still lags behind the command.

The intruder block combines current relative position with pre-computed
CPA quantities.
Relative position $(\tilde{x}_{fj},\,\tilde{y}_{fj})$ places the intruder
in the ego frame and is normalised by $d_\text{warn}$ to cover the full
conflict warning horizon.
The CPA coordinate pair $(\tilde{x}^\text{CPA}_{fj},\,\tilde{y}^\text{CPA}_{fj})$
encodes both the severity of the predicted conflict and, crucially, which
side the intruder will pass: $\tilde{x}^\text{CPA}_{fj} < 0$ indicates a
left-side conflict, providing a direct linear signal for turning direction.
These coordinates are normalised by $d_\text{sep}$ rather than $d_\text{warn}$
because the conflict geometry is confined to within one separation standard;
using $d_\text{warn}$ would compress the CPA features to a tiny range near
zero.
The normalised time $\hat{t}_\text{CPA}$ mirrors the imminence factor in
the reward function.
A value of 1 signals that no conflict is predicted, making the
CPA coordinate features uninformative (both zero).
Relative velocity is not included explicitly: it is implicit in the
combination of current position and CPA offset, from which the network can
recover the encounter geometry.
The CPA offset is defined as
\begin{equation}
    \begin{pmatrix}\Delta x^\text{CPA}_{fj}\\\Delta y^\text{CPA}_{fj}\end{pmatrix}
    =
    \begin{pmatrix}\Delta x_{fj}\\\Delta y_{fj}\end{pmatrix}
    + t_\text{CPA}
    \begin{pmatrix}\Delta v_{e,fj}\\\Delta v_{n,fj}\end{pmatrix},
\end{equation}
rotated into $f$'s heading frame before normalisation.
Non-converging pairs and empty slots receive the sentinel
$(1,\,1,\,0,\,1,\,1)$.

\subsection{Intruder Slot Ordering}
\label{sec:nbr}

The four intruder slots are filled in two passes.
In the first pass, all aircraft $j$ for which $u_{fj} > 0$ are inserted
in descending urgency order until the slots are full.
In the second pass, any remaining slots are filled by the nearest aircraft
by current distance, excluding those already selected.
This two-pass strategy ensures that predicted-conflict aircraft are
prioritised while nearby diverging aircraft, which may become relevant
within the current step, are guaranteed a slot when space permits.

% ─────────────────────────────────────────────────────────────────────────────
\section{Reward Function}
\label{sec:reward}
% ─────────────────────────────────────────────────────────────────────────────

The reward is purely negative at every step:

\begin{equation}
    r_t \;=\; r_\text{LoS} + r_\text{conflict} + r_\text{drift} + r_\text{work}
    \;\leq\; 0.
    \label{eq:reward}
\end{equation}

By excluding positive terms, the training objective reduces to minimising
cumulative harm, and the optimal policy is the one that keeps all penalties
as small as possible throughout the episode.
The four components and their weights are summarised in
Table~\ref{tab:weights}.

\paragraph{Loss-of-separation penalty.}
A separation violation is detected when any pair of aircraft is found within
$d_\text{sep}$ of each other at the end of the 5-second simulation window:

\begin{equation}
    r_\text{LoS} = -w_\text{LoS} \cdot \mathbf{1}[\ell_\text{LoS}].
\end{equation}

This component carries the largest weight and signals that a violation
has already occurred, as distinct from the graduated conflict signal below.
The check is sector-wide: violations between any pair, not only those
involving the focus aircraft, contribute to $\ell_\text{LoS}$.

\paragraph{Conflict penalty.}
For each converging intruder $k$ whose predicted closest-point-of-approach
satisfies $d_\text{CPA}^{(k)} < d_\text{sep}$, a severity score is
computed as the product of an imminence factor and a miss-distance factor:

\begin{equation}
    s_k \;=\; \max\!\left(0,\;1 - \frac{t_\text{CPA}^{(k)}}{T_\text{warn}}\right)
             \times
             \max\!\left(0,\;1 - \frac{d_\text{CPA}^{(k)}}{d_\text{sep}}\right),
    \label{eq:score}
\end{equation}

and the penalty is taken over the worst-case intruder:

\begin{equation}
    r_\text{conflict} = -w_\text{conflict} \cdot \max_k\, s_k.
\end{equation}

The imminence factor is zero when $t_\text{CPA} \geq T_\text{warn}$ and rises
to 1 as the conflict becomes imminent.
The miss-distance factor is zero when $d_\text{CPA} = d_\text{sep}$ and rises
to 1 as the predicted miss distance approaches zero, so that a grazing
predicted conflict scores lower than a direct collision course at the same
time horizon.
The predicted-LoS gate ($d_\text{CPA} < d_\text{sep}$) ensures that only
pairs on a true collision course contribute.
Aircraft already in active separation violation contribute $s_k = 1$.

\paragraph{Drift penalty.}
A continuous penalty on the commanded-heading deviation from the destination
bearing penalises the agent for holding the aircraft off its planned route:

\begin{equation}
    r_\text{drift}
    = -w_\text{drift} \cdot \frac{1 - \cos(\psi_\text{dest} - \psi_\text{cmd})}{2}.
\end{equation}

Unlike the workload term, drift is charged every step that the commanded
heading deviates from the destination bearing, regardless of which action
was issued.
During the hold phase this penalty accumulates, creating a sustained
incentive to issue the fly-direct action once a conflict has been resolved.
The penalty is zero when $\psi_\text{cmd} = \psi_\text{dest}$ and reaches
at most $(1 - \cos 60^\circ)/2 = 0.25$ at the largest deviation.

\paragraph{Workload penalty.}
A one-time cost is charged each time a turn instruction or a fly-direct
instruction is issued.
For turn instructions the cost is fixed:

\begin{equation}
    r_\text{work} = -w_\text{work} \cdot c,
\end{equation}

where $c$ is the action cost from Table~\ref{tab:actions}.
The hold action (action~3) is the only truly free action ($c = 0$).
For the fly-direct action (action~7) the cost scales with the commanded
heading deviation being corrected at the moment the instruction is issued:

\begin{equation}
    r_\text{work}^\text{direct}
    = -w_\text{work} \cdot \frac{|\Delta\psi_\text{cmd}|}{60^\circ},
    \qquad
    \Delta\psi_\text{cmd} = \psi_\text{dest} - \psi_\text{cmd}^{-},
\end{equation}

where $\psi_\text{cmd}^{-}$ denotes the commanded heading immediately before
the fly-direct is applied.
The cost is therefore zero when the aircraft is already on track and rises
to $w_\text{work}$ when correcting a full $60^\circ$ deviation, matching
the cost of the largest turn instruction.
This discourages the agent from prematurely aborting a conflict resolution
turn before the separation has been established, while keeping fly-direct
effectively free once the deviation is small.

\begin{table}[h]
\centering
\caption{Reward weights in descending order of magnitude.}
\label{tab:weights}
\smallskip
\begin{tabular}{clc}
\toprule
\textbf{Weight} & \textbf{Penalises} & \textbf{Value} \\
\midrule
$w_\text{LoS}$      & Separation violation at step end                                        & $10.00$ \\
$w_\text{conflict}$ & Worst-case imminence $\times$ miss-distance severity ($d_\text{CPA} < d_\text{sep}$) & $3.00$ \\
$w_\text{drift}$    & Per-step deviation; accumulates during hold, motivates fly-direct        & $0.60$ \\
$w_\text{work}$     & Per-turn cost (fixed); fly-direct cost $= w_\text{work}\cdot|\Delta\psi|/60^\circ$; hold free & $0.50$ \\
\bottomrule
\end{tabular}
\end{table}

% ─────────────────────────────────────────────────────────────────────────────
\section{Training Configuration}
\label{sec:training}
% ─────────────────────────────────────────────────────────────────────────────

The policy is trained with Proximal Policy Optimisation (PPO) as implemented
in Stable-Baselines3, using four parallel environments (\texttt{SubprocVecEnv})
each seeded independently.
Observation normalisation is disabled (\texttt{norm\_obs=False}) because all
features are bounded by design.
Reward normalisation is disabled (\texttt{norm\_reward=False}): the raw
reward signal is passed to the value function without scaling, so the
penalty magnitudes in Table~\ref{tab:weights} are the values seen during
training without transformation.
The policy and value networks share a two-layer MLP with $[128,\,128]$ hidden
units.

\begin{table}[h]
\centering
\caption{PPO hyperparameters.}
\label{tab:hp}
\smallskip
\begin{tabular}{llc}
\toprule
\textbf{Symbol} & \textbf{Parameter} & \textbf{Value} \\
\midrule
$\gamma$               & Discount factor                    & $0.99$           \\
$\lambda$              & GAE parameter                      & $0.95$           \\
$T$                    & Rollout steps per environment      & $1\,024$         \\
$N_\text{envs}$        & Parallel environments              & $4$              \\
$B$                    & Mini-batch size                    & $64$             \\
$K$                    & Epochs per update                  & $10$             \\
$\epsilon_\text{clip}$ & PPO clipping range                 & $0.2$            \\
$\beta_H$              & Entropy coefficient                & $0.02$           \\
$\alpha$               & Learning rate                      & $3\times10^{-4}$ \\
                       & Network architecture               & $[128,\,128]$    \\
\bottomrule
\end{tabular}
\end{table}

% ─────────────────────────────────────────────────────────────────────────────
\section{Validation}
\label{sec:validation}
% ─────────────────────────────────────────────────────────────────────────────

Policy performance is evaluated by sweeping a two-dimensional grid of
operating conditions: sector density $\rho$ on one axis and traffic count
$N$ on the other.
The sweep parameters are summarised in Table~\ref{tab:sweep}.
These ranges cover the high-density, high-count conditions that are most
relevant to operational stress testing, and extend beyond the training
distribution ($N \sim \mathcal{U}\{10,\ldots,15\}$) to assess
out-of-distribution robustness at both ends.

\begin{table}[h]
\centering
\caption{Validation sweep parameters.}
\label{tab:sweep}
\smallskip
\begin{tabular}{ll}
\toprule
\textbf{Parameter} & \textbf{Values} \\
\midrule
Sector density $\rho$ (km$^2$/aircraft) & $5\,000 \quad 8\,000 \quad 10\,000 \quad 15\,000$ \\
Traffic count $N$ (aircraft)            & $8 \quad 10 \quad 12 \quad 14$ \\
Episodes per cell                       & $8$ \\
Checkpoint                              & $4.6\,\text{M}$ training steps \\
\bottomrule
\end{tabular}
\end{table}

The primary metric is the \textit{loss-of-separation fraction}: the fraction
of steps within an episode in which at least one separation violation ($d < d_\text{sep}$)
is active.
Averaging over eight episodes per cell reduces variance from random sector
geometry and initial placement.

The same grid is evaluated under a \textit{no-policy baseline} in which
actions are drawn uniformly at random.
This baseline measures the separation violation rate that would occur without
any conflict resolution, so that the reduction attributable to the learned
policy can be quantified directly.

Results are presented as side-by-side heatmaps of mean loss-of-separation
fraction for the policy and the baseline, using a shared colour scale to
allow direct visual comparison.
A second figure shows loss-of-separation fraction as a function of density,
with one panel per traffic count and shaded bands indicating $\pm 1$ standard
deviation across episodes.

\end{document}
